% Lines 2045-2076 and 5857-5913 from original attention_transformers_notes_ghost.tex
% Appendix: Common Pitfalls and Solutions

%==============================================================================
\section{Common Pitfalls and Solutions}
%==============================================================================

\subsection{Training Issues}

\begin{table}[H]
\centering
\caption{Common Training Problems and Solutions}
\begin{tabular}{l l}
\toprule
\textbf{Problem} & \textbf{Solution} \\
\midrule
Gradient explosion & Gradient clipping, careful initialization \\
Training instability & Pre-norm architecture, warmup, smaller learning rate \\
Overfitting & Dropout, weight decay, data augmentation \\
Slow convergence & Learning rate scheduling, larger batch size \\
Memory overflow & Gradient accumulation, model parallelism \\
\bottomrule
\end{tabular}
\end{table}

\subsection{Inference Optimization}

\begin{enumerate}
    \item \textbf{KV-Cache}: Cache key-value pairs for autoregressive generation
    \item \textbf{Beam Search Optimization}: Batch beam search computations
    \item \textbf{Quantization}: INT8 quantization for deployment
    \item \textbf{Pruning}: Remove less important attention heads
    \item \textbf{Knowledge Distillation}: Train smaller models from larger ones
\end{enumerate}

% --- Continuation from line 5857 ---

%==============================================================================
\section{Common Pitfalls and Solutions}

\subsection{Training Issues}

\textbf{Problem: Gradient Explosion}
\begin{itemize}
    \item \textbf{Symptoms:} Loss becomes NaN, weights explode
    \item \textbf{Solutions:}
        - Implement gradient clipping: \texttt{torch.nn.utils.clip\_grad\_norm\_}
        - Use pre-layer normalization
        - Reduce learning rate
        - Check for numerical instabilities in attention
\end{itemize}

\textbf{Problem: Slow Convergence}
\begin{itemize}
    \item \textbf{Symptoms:} Loss plateaus early, poor performance
    \item \textbf{Solutions:}
        - Implement proper learning rate warmup
        - Increase model capacity (depth/width)
        - Check data preprocessing and tokenization
        - Verify positional encodings are working
\end{itemize}

\textbf{Problem: Out of Memory}
\begin{itemize}
    \item \textbf{Symptoms:} CUDA OOM errors
    \item \textbf{Solutions:}
        - Implement gradient accumulation
        - Use gradient checkpointing
        - Reduce batch size or sequence length
        - Consider efficient attention variants
\end{itemize}

\subsection{Implementation Bugs}

\textbf{Problem: Attention Mask Errors}
\begin{itemize}
    \item \textbf{Symptoms:} Model attends to padding or future tokens
    \item \textbf{Solutions:}
        - Verify mask dimensions match attention scores
        - Ensure masks use correct boolean values
        - Check mask is applied before softmax
\end{itemize}

\textbf{Problem: Dimension Mismatches}
\begin{itemize}
    \item \textbf{Symptoms:} Runtime errors in matrix operations
    \item \textbf{Solutions:}
        - Carefully track batch, sequence, and feature dimensions
        - Use einsum for complex operations
        - Add shape assertions throughout forward pass
\end{itemize}
